\begin{abstract}
DiCo-NLI asks systems to label ordered phrase pairs as equivalence, forward entailment, backward entailment or another relation, and scores whether predictions for the two directions of a pair are compatible under the reversal map. Rather than hoping that a classifier learns this algebra, we enforce it at inference: a multilingual encoder scores both directions, and \emph{joint bidirectional decoding} selects the best label for the pair as a whole, which makes the two predictions consistent by construction. Across English, Spanish, Basque and mixed-language tracks, joint decoding improves paired accuracy (HardCons) in all 72 of our evaluations, with no extra training, and it is more effective than a training-time consistency loss. Our best systems, 3-seed ensembles of NLI-pretrained or multilingually pooled mDeBERTa with joint decoding, reach about 79 weighted F1 and 83 HardCons on average on the dev sets. Error analysis shows that, for fine-tuned encoders, inconsistent directions are a minor error source; the main one is distinguishing reversible from non-reversible relations. We also find that SoftCons is maximised by a constant predictor, and that non-reversible items can be identified from the pair structure of the data alone.
\end{abstract}
